\documentclass[11pt]{article}

\usepackage[letterpaper,margin=0.82in,headheight=15pt]{geometry}
\usepackage[T1]{fontenc}
\usepackage{lmodern}
\usepackage{microtype}
\usepackage{amsmath,amssymb,bm,mathtools}
\usepackage{booktabs,longtable,tabularx,array}
\usepackage{enumitem}
\usepackage{xcolor}
\usepackage{fancyhdr}
\usepackage{hyperref}
\usepackage[nameinlink,noabbrev]{cleveref}
\usepackage{graphicx}
\usepackage{caption}

\definecolor{navy}{HTML}{17324D}
\definecolor{teal}{HTML}{147D92}
\definecolor{orange}{HTML}{D47A27}
\definecolor{softblue}{HTML}{EEF4F7}
\definecolor{softgray}{HTML}{F3F5F6}
\definecolor{darkgray}{HTML}{46545E}

\hypersetup{
  colorlinks=true,
  linkcolor=navy,
  citecolor=teal,
  urlcolor=teal,
  pdftitle={Construction and status of the canonical spin-1 deuteron GTMD/TMD model},
  pdfauthor={DeuteronWigner project}
}

\pagestyle{fancy}
\fancyhf{}
\fancyhead[L]{\small\color{darkgray}Canonical spin-1 deuteron GTMD/TMD model}
\fancyhead[R]{\small\color{darkgray}Pre-evolution technical note}
\fancyfoot[L]{\small\color{darkgray}27 July 2026}
\fancyfoot[R]{\small\color{darkgray}\thepage}
\renewcommand{\headrulewidth}{0.4pt}
\renewcommand{\footrulewidth}{0pt}

\setlist[itemize]{leftmargin=1.5em,itemsep=2pt,topsep=3pt}
\setlist[enumerate]{leftmargin=1.7em,itemsep=2pt,topsep=3pt}
\setlength{\parindent}{0pt}
\setlength{\parskip}{5pt}
\setlength{\emergencystretch}{2em}
\renewcommand{\arraystretch}{1.15}

\newcommand{\kT}{\bm{k}_{T}}
\newcommand{\pT}{\bm{p}_{T}}
\newcommand{\bT}{\bm{b}_{T}}
\newcommand{\DT}{\bm{\Delta}_{T}}
\newcommand{\RT}{\bm{R}_{T}}
\newcommand{\Tr}{\operatorname{Tr}}
\newcommand{\dd}{\mathrm{d}}
\newcommand{\LF}{\mathrm{LF}}
\newcommand{\GTMD}{\mathrm{GTMD}}
\newcommand{\TMD}{\mathrm{TMD}}
\newcommand{\CS}{\mathrm{CS}}
\newcommand{\IA}{\mathrm{IA}}
\newcommand{\status}[1]{\textsf{\textbf{#1}}}
\newcommand{\exact}{\status{Exact}}
\newcommand{\fit}{\status{Fit/phenomenology}}
\newcommand{\model}{\status{Model}}
\newcommand{\temporary}{\status{Temporary}}
\newcommand{\implemented}{\textcolor{teal}{\status{Implemented}}}
\newcommand{\openitem}{\textcolor{orange}{\status{Open}}}

\newcolumntype{Y}{>{\raggedright\arraybackslash}X}
\newcolumntype{L}[1]{>{\raggedright\arraybackslash}p{#1}}

\title{\vspace{-1.0cm}
  {\color{navy}\bfseries Construction and Status of the Canonical\\[2mm]
  Spin-1 Deuteron GTMD/TMD Model}\\[4mm]
  \large A light-front, flavor-resolved, spin-resolved, quark--gluon boundary
  before complete rank-aware evolution}
\author{DeuteronWigner project}
\date{27 July 2026}

\begin{document}
\maketitle

\begin{center}
\begin{minipage}{0.93\textwidth}
\colorbox{softblue}{\parbox{0.96\textwidth}{
\textbf{Scope statement.}
This is the authoritative scientific and implementation note for the
leading-twist forward boundary at
\(Q=5~\mathrm{GeV}\) and
\(x_N\in\{0.02,0.05,0.10,0.20,0.40\}\).
The model contains all declared spin-1 quark, antiquark, and gluon
projections, retains proton and neutron parents, and separates exact
constraints, fitted inputs, model assumptions, and unresolved physics.
Complete common-scheme multi-\(Q\) evolution is the next phase and is not
claimed here.
}}
\end{minipage}
\end{center}

\begin{abstract}
We describe the construction, physical content, validation, and present
status of a canonical light-front model for the complete declared
leading-twist spin-1 deuteron transverse-momentum-dependent (TMD) system.
The organizing object is a retained quark or gluon parton--target helicity
correlator obtained from a nucleon correlator and a spin-dependent
light-front deuteron spectral density.  Named TMDs are projections of this
parent, rather than independent curves.  The boundary is explicitly
resolved in \(u,d,\bar u,\bar d\), proton and neutron source, target
polarization \(U,L,T,LL,LT,TT\), orbital interference, nuclear mechanism,
and, for naive-T-odd gluons, the antisymmetric \(f^{abc}\) and symmetric
\(d^{abc}\) gauge-link/color classes.

The numerical construction combines CT18 and MSHT20QED unpolarized inputs,
BDSSV24 helicity replicas, JAMDiFF plus lattice-informed transversity,
BPV20 Sivers replicas, a published-interval treatment of the Yang
\(g_{1T}\) fit, physically constrained Boer--Mulders, pretzelosity,
worm-gear, tensor, and gluon sectors, six realistic deuteron wave
functions, binding and Fermi motion, off-shell response, shadowing and
antishadowing, a sourced \(NN\pi\) component, and controlled non-nucleonic
interfaces.  We give the field-theoretic definitions, polarization and
transverse-rank algebra, nuclear convolution, composition rules,
uncertainty semantics, and validation conditions.  We also document why
two earlier implementations---a real one-body boundary and a smooth
reduced-amplitude closure---were insufficient.  The accepted result is a
physically organized pre-evolution boundary, not a global extraction and
not a claim of equal evidential precision for all TMDs.
\end{abstract}

\tableofcontents
\clearpage

\section{Scientific objective and status language}
\label{sec:objective}

The objective is not a minimally runnable atlas.  It is a model in which
all realistically supported contributions to the spin-1 partonic density
can coexist without losing their origin or being counted twice.  A model
is called \emph{canonical} here when
\begin{enumerate}
  \item the dynamical state is a retained quark--target or gluon--target
  correlator;
  \item every named TMD is obtained by a declared projection of that
  correlator;
  \item proton, neutron, flavor, helicity, wave-function component,
  gauge-link, color, and nuclear-mechanism identity remain available until
  the physical observable is assembled;
  \item a machine-readable composition graph distinguishes baselines,
  additive mechanisms, mutually exclusive alternatives, and ensemble
  members;
  \item exact constraints, fit information, lattice information, and model
  assumptions remain different provenance classes.
\end{enumerate}

The word \emph{complete} must always carry a scope.  In this note it means:
the complete declared leading-twist forward quark and gluon basis at the
fixed boundary scale \(Q=5~\mathrm{GeV}\), with smooth functions and named
bands over the five production \(x_N\) nodes.  It does not mean a
first-principles solution of QCD, a simultaneous global fit, complete
nonzero-skewness GTMD phenomenology, or complete TMD evolution.

\subsection{Evidence vocabulary}

We use the following classification throughout:
\begin{description}[leftmargin=2.2cm,style=nextline]
  \item[\exact] A consequence of the chosen operator definition,
  representation, discrete symmetry, kinematics, or algebraic composition
  rule.
  \item[\fit] A central function or ensemble tied to a published PDF/TMD
  extraction or directly calibrated observable.
  \item[\status{Lattice-informed}] A fit or model component constrained by
  lattice matrix elements or an explicit lattice prior.
  \item[\model] A physically motivated but underdetermined component with a
  replaceable parameterization and a named sensitivity ensemble.
  \item[\temporary] A controlled implementation boundary retained only
  until a specified replacement, with a test that distinguishes the two.
\end{description}
Passing a structural acceptance gate does not upgrade a model prediction
to a fit.

\section{From the original proposal to the present model}
\label{sec:history}

\subsection{The GTMD-first proposal}

The original technical draft, \emph{A GTMD-First Light-Front Wigner and
SCET Framework for Spin-1 Deuteron TMDs}, proposed the correct parent
chain, following the GTMD/Wigner logic of
Refs.~\cite{Belitsky2004,Lorce2011}:
\begin{equation}
\Psi_D^{\LF}
\longrightarrow
\rho_{\Lambda'\Lambda;\lambda'\lambda}^{N/D}
\longrightarrow
W_{a/N}^{[\Gamma]}
\longrightarrow
W_{a/D}^{[\Gamma]}
\longrightarrow
\Phi_{a/D}^{[\Gamma]}
\longrightarrow
F_{a/D}.
\label{eq:parent-chain}
\end{equation}
The light-front deuteron wave function \(\Psi_D^{\LF}\) determines the
active-nucleon spin density \(\rho^{N/D}\); this density convolves with a
nucleon partonic GTMD \(W_{a/N}\); the forward correlator
\(\Phi_{a/D}\) is projected onto named functions \(F_{a/D}\).

The parent structure is important because the reductions must commute:
\begin{align}
F_{a/D}(x,\kT;\eta)
  &= W_{a/D}(x,\kT,\DT=\bm 0;\eta),
  &&\text{TMD limit}, \label{eq:tmd-limit}\\
H_{a/D}(x,\DT)
  &= \int \dd^2\kT\,W_{a/D}(x,\kT,\DT),
  &&\text{zero-skewness GPD limit}, \label{eq:gpd-limit}\\
f_{a/D}(x)
  &= \int \dd^2\kT\,F_{a/D}(x,\kT),
  &&\text{rank-zero PDF limit}, \label{eq:pdf-limit}\\
\mathcal W_{a/D}(x,\kT,\bm b_\Delta)
  &= \int\frac{\dd^2\DT}{(2\pi)^2}
       e^{-i\bm b_\Delta\cdot\DT}W_{a/D}(x,\kT,\DT),
  &&\text{partonic Wigner transform}.
\label{eq:wigner}
\end{align}
The nuclear Wigner/spectral density and the partonic Wigner distribution
are distinct: the first describes nucleons in the deuteron; the second is
the transverse Fourier transform of a QCD partonic correlator.

\subsection{First refocusing: exact zeros were boundary zeros}

The first numerical parent used a real, helicity-independent, rank-zero
nucleon GTMD.  It correctly tested Fourier signs, recoil, hermiticity,
marginals, and the tensor PDF entering \(b_1\).  It could not represent the
complete spin structure.  If all amplitudes are real,
\begin{equation}
\operatorname{Im}
\left[\psi_{L_z'}^\ast\,\mathcal U_\eta\,\psi_{L_z}\right]=0
\quad\Longrightarrow\quad
F_{\text{naive-T-odd}}=0,
\label{eq:real-zero}
\end{equation}
and missing off-diagonal helicity or OAM amplitudes similarly force
worm-gear, pretzelosity, and tensor projections to vanish.  These were
controlled one-body limits, not predictions of exact physical zeros.

\subsection{Second refocusing: smooth closure was not parent dynamics}

A later reduced-amplitude construction populated every TMD name with a
smooth curve and conservative rank factors.  Its map was schematically
\begin{equation}
\{f_1,g_1,h_1,f_{1LL}\}_{\rm anchors}
\longrightarrow
\{\mathcal A_U,\mathcal A_L,\mathcal A_T,
\mathcal A_{LL},\mathcal A_{LT},\mathcal A_{TT}\}
\longrightarrow
\{F_A\}_{\rm all\ names}.
\label{eq:bad-closure}
\end{equation}
This was useful for registry and plotting tests but bypassed
\cref{eq:parent-chain}.  Shared shapes and coefficients produced
flavor-independent appearances, while proton and neutron sources were
summed too early.  Passing smoothness and individual positivity ceilings
did not establish a physical model.  Those outputs are now explicitly
superseded.

\subsection{Third refocusing: isoscalar equality belongs at final assembly}

In the exact charge-symmetric one-body limit, an inclusive \(I=0\)
deuteron satisfies
\begin{equation}
u_D=u_{p/D}+u_{n/D}=u_{p/D}+d_{p/D}=d_D,
\qquad
\bar u_D=\bar d_D.
\label{eq:isoscalar}
\end{equation}
This does not imply \(u_p=d_p\), nor does it justify erasing constituent
identity.  The production state therefore retains
\begin{multline*}
\{\text{proton-in-deuteron},\ \text{neutron-in-deuteron},\
\text{nucleon sum},\\
\text{proton-minus-neutron},\ \text{nuclear correction},\
\text{spin-1 total}\}.
\end{multline*}
The spin-1 total is a closure projection, not the stored model state.

\subsection{Fourth refocusing: basis parity was not evidence parity}

After the correlator basis was populated, a final audit required every one
of the 36 declared quark-plus-gluon projections to identify a central
source, constituent resolution, uncertainty or sensitivity ensemble, CSB
status, shared-parent projection, channel-appropriate nuclear dressing, and
a physical benchmark or controlled limit.  All 36 rows now pass this
evidence-discipline gate.  Many tensor and gluon rows remain predictions,
not direct extractions.

\section{Kinematics and spin-1 representation}
\label{sec:representation}

We use lightlike vectors \(n^\mu,\bar n^\mu\) with
\(n^2=\bar n^2=0\) and \(n\cdot\bar n=1\):
\begin{equation}
v^\mu=v^+\bar n^\mu+v^-n^\mu+v_T^\mu,\qquad
v^+=v\cdot n,\quad v^-=v\cdot\bar n.
\end{equation}
At zero skewness the average deuteron momentum is \(P\), the transverse
transfer is \(\DT\), and the active parton has average transverse momentum
\(\kT\).  The staple direction is denoted \(\eta=+\) (future) or
\(\eta=-\) (past).  Nuclear convolutions use the nucleon-scaled fraction
\(x_N\), with \(x_D=x_N/2\) in the project convention.

\subsection{Target density and irreducible basis}

The target Hilbert space is
\(\mathcal H_D=\operatorname{span}\{|+1\rangle,|0\rangle,|-1\rangle\}\).
Every Hermitian target matrix has the irreducible decomposition
\begin{equation}
\operatorname{Herm}(3)
=
\underbrace{U}_{1}
\oplus\underbrace{L}_{1}
\oplus\underbrace{(T_x,T_y)}_{2}
\oplus\underbrace{LL}_{1}
\oplus\underbrace{(LT_x,LT_y)}_{2}
\oplus\underbrace{(TT_x,TT_y)}_{2}.
\label{eq:spin1-decomp}
\end{equation}
The dimensions sum to nine.  With trace-orthogonal basis matrices
\(B_A\), a target matrix \(M\) is reconstructed as
\begin{equation}
M=\sum_A c_A B_A,\qquad
c_A=\sum_B(G^{-1})_{AB}\Tr(B_B^\dagger M),\qquad
G_{AB}=\Tr(B_A^\dagger B_B).
\label{eq:gram-project}
\end{equation}
This is the concrete representation-theoretic machinery used in the code.
It prevents incompatible target channels from being mixed and exposes
unidentifiable combinations through the Gram rank.

For helicity-resolved scalar quantities \(F_\Lambda\), useful combinations
are
\begin{align}
F_U&=\frac13(F_++F_0+F_-),&
F_L&=\frac12(F_+-F_-),&
\delta_TF&=F_0-\frac12(F_++F_-).
\label{eq:helicity-combos}
\end{align}
The internal \(LL\) basis has the opposite sign to the physical Trento
\(S_{LL}\) convention; explicit adapters are therefore required for
\(f_{1LL}\) and \(h_{1LL}^{\perp}\).

\subsection{Transverse rank}

The transverse plane carries \(SO(2)\).  A rank-\(r\) coefficient
\(F^{(r)}\) multiplies an irreducible harmonic
\(k_T^{i_1\cdots i_r}\), and its physical modulation scales as
\begin{equation}
\mathcal M_r(x,\kT)\sim
\left(\frac{k_T}{M_D}\right)^rF^{(r)}(x,k_T^2).
\label{eq:rank-weight}
\end{equation}
A positive-rank coefficient can be finite at \(k_T=0\), while the physical
modulation must vanish.  Rank determines the Fourier--Bessel kernel:
\begin{align}
\widetilde F^{(r)}(x,b_T)
 &=2\pi\int_0^\infty k_T\,\dd k_T\,
 J_r(b_Tk_T)\,\mathcal N_r(k_T)F^{(r)}(x,k_T),\\
F^{(r)}(x,k_T)
 &=\int_0^\infty\frac{b_T\,\dd b_T}{2\pi}\,
 J_r(b_Tk_T)\,\widetilde{\mathcal N}_r(b_T)
 \widetilde F^{(r)}(x,b_T).
\label{eq:bessel}
\end{align}
The registry stores \(r\) and the mass convention so that a scalar
\(J_0\) transform cannot be applied to rank-one or rank-two functions.

\section{Operator correlators and gauge links}
\label{sec:operators}

\subsection{Quark and antiquark correlators}

For flavor \(q\), the unsubtracted zero-skewness quark GTMD correlator is
\begin{equation}
W_{\Lambda'\Lambda}^{q[\Gamma]}(x,\kT,\DT;\eta)
=\frac12\int\frac{\dd z^-\,\dd^2\bm z_T}{(2\pi)^3}
e^{ixP^+z^- -i\kT\cdot\bm z_T}
\left.
\langle P',\Lambda'|
\bar q(-z/2)\,\Gamma\,\mathcal U_\eta[-z/2,z/2]\,
q(z/2)
|P,\Lambda\rangle
\right|_{z^+=0}.
\label{eq:quark-gtmd}
\end{equation}
The leading-twist projections are
\(\Gamma=\gamma^+\), \(\gamma^+\gamma_5\), and
\(i\sigma^{i+}\gamma_5\)~\cite{BacchettaMulders2000,KumanoSong2021}.
The physical TMD is UV-renormalized and
soft-subtracted.  Antiquarks are represented by a separate positive-\(x\)
sector rather than by silently identifying them with quarks.

At fixed \((x,\kT)\), the retained quark object is a \(3\times3\) target
matrix for each Dirac projection, equivalently a \(6\times6\) joint
target--quark helicity density when the quark helicity indices are made
explicit.

\subsection{Gluon correlator}

The gluon correlator and its gauge-link labels follow the general spin
\(\leq1\) classification of Ref.~\cite{Boer2016}:
\begin{align}
\Gamma_{\Lambda'\Lambda}^{ij}(x,\kT,\DT;\eta_1,\eta_2)
&=\frac{1}{xP^+}\int\frac{\dd z^-\,\dd^2\bm z_T}{(2\pi)^3}
e^{ixP^+z^- -i\kT\cdot\bm z_T}
\nonumber\\[-1mm]
&\quad\times
\left.
\langle P',\Lambda'|
\Tr_c\!\left[
F^{+i}(-z/2)\mathcal U_{\eta_1}
F^{+j}(z/2)\mathcal U_{\eta_2}
\right]|P,\Lambda\rangle
\right|_{z^+=0}.
\label{eq:gluon-gtmd}
\end{align}
The transverse \(2\times2\) gluon space decomposes into trace
(unpolarized), antisymmetric (circular polarization), and
symmetric-traceless (linear polarization) pieces.  Together with the
target space, the stored forward parent is a
\(3\times3\times2\times2\) joint matrix.

\subsection{Discrete symmetries and process dependence}

Hermiticity requires
\begin{equation}
W_{\Lambda'\Lambda}(x,\kT,\DT;\eta)
=W_{\Lambda\Lambda'}^\ast(x,\kT,-\DT;\eta).
\label{eq:hermiticity}
\end{equation}
Parity is implemented as a tested light-front helicity reflection.  Time
reversal maps future to past gauge links:
\begin{equation}
F_{\rm even}^{[+]}=F_{\rm even}^{[-]},\qquad
F_{\rm odd}^{[+]}=-F_{\rm odd}^{[-]}.
\label{eq:link-reversal}
\end{equation}
For gluons, the two universal naive-T-odd color contractions
\begin{equation}
\mathcal C_f^{abc}=if^{abc},\qquad
\mathcal C_d^{abc}=d^{abc}
\label{eq:fd}
\end{equation}
remain independent.  A process observable has the form
\begin{equation}
F_{\rm process}^g=C_f^{\rm process}F^{g(f)}
                 +C_d^{\rm process}F^{g(d)}.
\label{eq:color-combination}
\end{equation}
The model refuses an implicit universal \(f/d\) mixture or a mixed-link
sign assigned by quark analogy.

\section{Complete declared leading-twist basis}
\label{sec:basis}

The current forward basis follows
Refs.~\cite{Spin1Quark2016,KumanoSong2021,Boer2016}.  It has 18 quark functions for each of
\(u,d,\bar u,\bar d\), and 18 identifiable gluon projections.  The older
exploratory 19-gluon count split a TT pair that the implemented forward
correlator identifies only through a combination; it is not the canonical
registry.

\begin{table}[ht]
\centering
\caption{Declared quark and antiquark basis for each flavor.  The rank is
the rank of the transverse tensor multiplying the named scalar coefficient.}
\label{tab:quark-basis}
\small
\begin{tabularx}{\textwidth}{L{0.10\textwidth}L{0.56\textwidth}Y}
\toprule
Target & Named functions & Physical content \\
\midrule
\(U\)  & \(f_1,\ h_1^\perp\) & unpolarized density; Boer--Mulders \\
\(L\)  & \(g_1,\ h_{1L}^\perp\) & helicity; longitudinal worm gear \\
\(T\)  & \(f_{1T}^\perp,\ g_{1T},\ h_1,\ h_{1T}^\perp\) &
Sivers, transverse worm gear, transversity, pretzelosity \\
\(LL\) & \(f_{1LL},\ h_{1LL}^\perp\) &
longitudinal tensor density and linear/chiral-odd partner \\
\(LT\) & \(f_{1LT},\ g_{1LT},\ h_{1LT},\ h_{1LT}^\perp\) &
mixed tensor polarization \\
\(TT\) & \(f_{1TT},\ g_{1TT},\ h_{1TT},\ h_{1TT}^\perp\) &
transverse tensor polarization \\
\bottomrule
\end{tabularx}
\end{table}

\begin{table}[ht]
\centering
\caption{Declared gluon basis.  The f-type and d-type gauge-link/color
sectors are additional labels on T-odd entries, not additional TMD names.}
\label{tab:gluon-basis}
\small
\begin{tabularx}{\textwidth}{L{0.10\textwidth}L{0.58\textwidth}Y}
\toprule
Target & Named functions & Physical content \\
\midrule
\(U\) & \(f_1^g,\ g_1^g,\ h_1^{\perp g}\) &
unpolarized, helicity, linear polarization \\
\(L\) & \(h_{1L}^{\perp g}\) & longitudinal spin--linear-polarization interference \\
\(T\) & \(f_{1T}^{\perp g},\ g_{1T}^g,\ h_1^g,\ h_{1T}^{\perp g}\) &
gluon Sivers, worm gear, and transverse spin/linear structures \\
\(LL\) & \(f_{1LL}^g,\ h_{1LL}^{\perp g}\) & longitudinal tensor gluons \\
\(LT\) & \(f_{1LT}^g,\ g_{1LT}^g,\ h_{1LT}^g,\ h_{1LT}^{\perp g}\) &
mixed tensor gluons \\
\(TT\) & \(f_{1TT}^g-h_{1TT}^{\perp g},\
g_{1TT}^g,\ h_{1TT}^g,\ h_{1TT}^{\perp\perp g}\) &
identifiable transverse-tensor combinations \\
\bottomrule
\end{tabularx}
\end{table}

\section{Nucleon-level quark dynamics}
\label{sec:quarks}

\subsection{Unpolarized and helicity sectors}

The central collinear \(f_1^q\) uses CT18NNLO~\cite{CT18} with distinct
\(u,d,\bar u,\bar d\) proton functions.  The 29 CT18 Hessian pairs are
propagated through the resolved nuclear response on all five production
\(x_N\) nodes.  The neutron is not obtained by assigning the same function
to both flavors.  Exact charge symmetry supplies the baseline rotation
\begin{equation}
u_n=d_p,\quad d_n=u_p,\quad
\bar u_n=\bar d_p,\quad \bar d_n=\bar u_p,
\label{eq:charge-symmetry}
\end{equation}
and MSHT20QED~\cite{MSHTQED} supplies a numerical unpolarized CSB ratio
\begin{equation}
\delta_q^n(x,Q)=
\frac{q_{\rm MSHT20QED}^{n}(x,Q)}
     {q_{\rm MSHT20QED}^{p,\,{\rm isospin\ partner}}(x,Q)}-1
\label{eq:csb-ratio}
\end{equation}
with 38 paired Hessian directions.  This numerical correction is not copied
to polarized or T-odd channels.

Quark \(g_1\) uses BDSSV24-NLO~\cite{BDSSV24}.  All 600 replicas are propagated through
the AV18 response and through the production x grid.  The wave-function
spin depolarization emerges from the helicity convolution; it is not
implemented only as a universal scalar.

\subsection{Transversity and chiral-odd worm gear}

The \(h_1\) input is the JAMDiFF \(w\)LQCD ensemble~\cite{JAMDiff}:
member zero plus 968
physical replicas.  The combination is fit and lattice informed.  Because
the source fit and the project \(f_1,g_1\) use different underlying PDF
sets, each member is projected into the actual TMD-level Soffer interval
\begin{equation}
|h_1^q(x,k_T)|\leq
\frac12\left[f_1^q(x,k_T)+g_1^q(x,k_T)\right]
\label{eq:soffer}
\end{equation}
with a small numerical interior margin.  Raw source members remain
recoverable, so this compatibility projection is visible and replaceable.

The \(h_{1L}^{\perp}\) boundary is correlated with the same replica identity
through a Wandzura--Wilczek-type functional.  Schematically,
\begin{equation}
h_{1L}^{\perp(1)q}(x)
\simeq -x^2\int_x^1\frac{\dd y}{y^2}\,h_1^q(y)
\;+\;\Delta_{\rm genuine}^q(x).
\label{eq:ww-h1l}
\end{equation}
The first term is an organized approximation; the genuine quark--gluon
term is an independent model sensitivity, not absorbed into the fit band.

\subsection{Sivers and Boer--Mulders}

The quark Sivers function uses the BPV20 N3LO arTeMiDe release
\cite{BPV20} with 500
replicas.  Member identity is propagated through all six nuclear wave
functions.  The source fit's light-sea equality is retained as a documented
fit assumption rather than generalized into a project-wide flavor
identity.  Staple reversal follows \cref{eq:link-reversal}.

Boer--Mulders \(h_1^\perp\) is represented by flavor-dependent
sign/hierarchy scenarios informed by the fitted Sivers rescattering
strength.  The relation
\begin{equation}
h_1^{\perp q}(x,k_T)
=C_{\rm BM}^q(x,k_T)\,f_{1T}^{\perp q}(x,k_T)
\label{eq:bm-model}
\end{equation}
is explicitly \model: \(C_{\rm BM}^q\) has alternative members and is not a
claim of an operator identity or a joint fit.

\subsection{Worm gear \(g_{1T}\) and pretzelosity}

The central \(g_{1T}\) parameters follow Yang et al.\ (2024)
\cite{Yang2024}.  Because no
public covariance or replica ensemble is available, the implementation
propagates all 16 corners of the four published asymmetric parameter
intervals.  The resulting hull is not labeled a confidence interval.  A
shared-Fock/OAM sea member tests the source fit's zero-sea boundary.

Pretzelosity is not generated by rescaling transversity; its unusual
small-\(b_T\) status is discussed in Ref.~\cite{Pretzel}.  Its
nonperturbative first transverse moment is parameterized inside
\begin{equation}
h_{1T}^{\perp(1)q}(x)
=\int\dd^2\kT\,\frac{k_T^2}{2M_N^2}
h_{1T}^{\perp q}(x,k_T^2),\qquad
\left|h_{1T}^{\perp(1)q}\right|
\leq\frac{f_1^q-g_1^q}{2}.
\label{eq:pretzel-bound}
\end{equation}
The central perturbative small-\(b_T\) boundary is zero, while signed
hadronic OAM members use a conservative fraction of
\cref{eq:pretzel-bound}.  These members are model sensitivities.

\subsection{Tensor-polarized quarks}

The ten LL/LT/TT quark functions are projections of the common deuteron
spin-density parent.  Their sources are S--D and OAM interference, Wilson
channels, spin recoupling, and operator-valued nuclear response.  In
particular, \(g_{1LT}\) and \(g_{1TT}\) require absorptive axial-tensor
interference and are implemented in two mutually exclusive stages:
\begin{enumerate}
  \item an independent positivity-bounded phase ensemble for each flavor
  and operator;
  \item a screened one-gluon eikonal model with rank \(n=1,2\):
\end{enumerate}
\begin{equation}
\mathcal I_n(k_T)=\frac{C_F\alpha_s}{2\pi}
\int_0^{q_{\max}}q\,\dd q\,
\frac{\left[\Lambda^2/(\Lambda^2+q^2)\right]^2}
     {q^2+\mu_g^2}
\left(\frac{q}{M_D}\right)^n
\left\langle
\cos(n\phi)e^{-|\kT-\bm q|^2/(2w)}
\right\rangle_\phi .
\label{eq:eikonal}
\end{equation}
The rank-one term represents imaginary S--P interference; the rank-two
term includes imaginary S--D and P-even/P-odd interference.  Disabling the
imaginary amplitude makes both functions vanish.  The larger direct-phase
envelope is not used to inflate the screened calculation.

\section{Nucleon-level gluon dynamics}
\label{sec:gluons}

\subsection{Unpolarized, helicity, and linearly polarized anchors}

The canonical \(f_1^g\) central boundary uses the matched
BSV19/NNPDF31 construction~\cite{BSV19}.  A separate CT18 29-pair Hessian response is
propagated through AV18 smearing at five x nodes and 61 transverse-momentum
points, then anchored as an uncertainty axis around the common central
scheme.  This preserves the distinction between the central TMD
construction and collinear-PDF uncertainty.

Gluon \(g_1^g\) uses all 600 BDSSV24-NLO replicas~\cite{BDSSV24}.
Width variations are
stored separately: they can strongly reshape the value near \(k_T=0\)
while leaving the finite-grid collinear integral nearly invariant.

The T-even \(h_1^{\perp g}\), \(g_{1T}^g\), \(h_1^g\), and
\(h_{1T}^{\perp g}\) arise from the retained gluon light-front
overlap/OAM/Wilson parent.  They are not quark TMDs with changed labels.
The gluon polarization tensor and the full joint density determine the
allowed interference and rank.

\subsection{Naive-T-odd gluons}

An early spectator-inspired model multiplied the assembled deuteron
\(f_1^g\) by literature-motivated coefficients.  It reproduced qualitative
hierarchies but bypassed the nucleon parent, and is benchmark-only.

The canonical construction instead introduces the absorptive Wilson
amplitude at the nucleon-correlator level, retains the two color classes in
\cref{eq:fd}, and then performs the same nuclear helicity convolution as
the T-even sector.  External spectator and CGI-GPM calculations constrain
sign, node, and hierarchy scenarios, but no released global replica set
determines the absolute scale of all gluon T-odd functions.  Process
predictions must therefore expose \(C_f,C_d\) in
\cref{eq:color-combination}.

\subsection{Tensor-polarized gluons}

The LL, LT, and TT gluon projections use the same shared
light-front/OAM/Wilson parent, projected with the spin-1 target basis.
The previously missing \(g_{1LT}^g\) and \(g_{1TT}^g\) are generated by
S--D/OAM interference and screened eikonal phases, with pure-S,
zero-phase, rank-covariance, link-reversal, and positivity tests.  Their
absolute magnitudes are predictions with named model axes.

\section{Light-front deuteron and nuclear composition}
\label{sec:nuclear}

\subsection{Two-nucleon light-front state}

The nucleonic component is
\begin{equation}
|D,\Lambda\rangle_{NN}
=\sum_{\lambda_1,\lambda_2}
\int\frac{\dd y\,\dd^2\pT}
          {2(2\pi)^3\,y(1-y)}
\Psi_{\Lambda;\lambda_1\lambda_2}^{D,\LF}(y,\pT)
|N_1(y,\pT,\lambda_1)N_2(1-y,-\pT,\lambda_2)\rangle .
\label{eq:deuteron-lf}
\end{equation}
The instant-form radial input contains
\begin{equation}
\Psi_\Lambda^{\rm IF}(\bm k)
=u_0(k)\,\mathcal Y_{00}^{1\Lambda}(\hat{\bm k})
u_2(k)\,\mathcal Y_{2}^{1\Lambda}(\hat{\bm k}),
\label{eq:sd-wave}
\end{equation}
followed by the declared light-front variable map and constituent Melosh
rotations.  The production ensemble contains AV18 (central)~\cite{AV18},
CD-Bonn~\cite{CDBonn}, and
Norfolk NV-Ia, NV-Ib, NV-IIa, and NV-IIb.

The active-nucleon density is
\begin{equation}
\rho_{\Lambda'\Lambda;\lambda'\lambda}^{N/D}(y,\pT)
=\sum_{\lambda_s}
\Psi_{\Lambda';\lambda'\lambda_s}^{D,\LF\,\ast}(y,\pT)
\Psi_{\Lambda;\lambda\lambda_s}^{D,\LF}(y,\pT),
\label{eq:spectral-density}
\end{equation}
with separate SS, SD, DS, and DD terms.  The D state is therefore not
represented only by one probability.  Signed S--D interference enters
tensor polarization and spin--orbit observables directly.

\subsection{GTMD and TMD convolution}

At zero skewness, the one-body nuclear GTMD has the schematic retained-spin
form
\begin{align}
W_{a/D,\Lambda'\Lambda}^{[\Gamma]}(x,\kT,\DT)
&=\sum_{N=p,n}
\int\frac{\dd y}{y}\frac{\dd^2\pT}{(2\pi)^2}
\sum_{\lambda',\lambda}
\rho_{\Lambda'\Lambda;\lambda'\lambda}^{N/D}
    (y,\pT,\DT)
\nonumber\\[-1mm]
&\quad\times
W_{a/N,\lambda'\lambda}^{[\Gamma]}
\left(\frac{x}{y},
\kT-\frac{x}{y}\pT,\bm{\Delta}_{T,N}\right)
J(y,\pT,\DT).
\label{eq:gtmd-convolution}
\end{align}
Its forward limit gives
\begin{equation}
F_{a/D,I}^{\IA}(x,\kT)
=\sum_{N=p,n}\int_x^1\frac{\dd y}{y}
\int\frac{\dd^2\pT}{(2\pi)^2}
\rho_{I;\lambda'\lambda}^{N/D}(y,\pT)
F_{a/N,\lambda'\lambda}
\left(\frac{x}{y},\kT-\frac{x}{y}\pT\right)J_T.
\label{eq:tmd-convolution}
\end{equation}
In \(b_T\) space the transverse shift becomes a phase:
\begin{align}
\widetilde F_{a/D,I}^{\IA}(x,\bT)
&=\sum_N\int_x^1\frac{\dd y}{y}
\int\frac{\dd^2\pT}{(2\pi)^2}
\rho_{I;\lambda'\lambda}^{N/D}(y,\pT)
e^{i(x/y)\bT\cdot\pT}
\widetilde F_{a/N,\lambda'\lambda}(x/y,\bT)J_T.
\label{eq:bspace-convolution}
\end{align}
This is the form to be consumed by complete evolution.

\subsection{Binding, off-shell response, and completely positive maps}

Binding and Fermi motion enter through \(y,\pT\) in
\cref{eq:tmd-convolution}.  Off-shell, shadowing, and antishadowing
responses act on the full joint spin density.  The accepted implementation
uses ordered completely positive maps
\begin{equation}
\mathcal E_m[\rho]
=\sum_\alpha K_{m\alpha}\rho K_{m\alpha}^\dagger,
\qquad
\sum_\alpha K_{m\alpha}^\dagger K_{m\alpha}\preceq \mathbf 1,
\label{eq:kraus}
\end{equation}
so Hermiticity and positivity are preserved.  The ordering is explicit
because polarized response maps need not commute.  These maps replace, not
add to, older scalar coefficient responses.

\subsection{Shadowing and antishadowing}

The unpolarized small-x baseline follows the H1 diffractive PDF
\cite{H1DPDF} and the Frankfurt--Guzey--Strikman double-scattering
construction~\cite{FGS}.  Schematically,
\begin{equation}
\delta f_{j/D}^{\rm shad}(x,Q)
\sim -\int_x^{x_{\mathbb P,\max}}\dd x_{\mathbb P}
\int\dd t\,
f_{\mathbb P/p}(x_{\mathbb P},t)\,
f_{j/\mathbb P}(x/x_{\mathbb P},Q)\,
F_D(4[q_T^2+x_{\mathbb P}^2m_N^2]).
\label{eq:shadowing}
\end{equation}
Polarized and tensor shadowing responses are independent model ratios:
unpolarized diffraction does not determine them.  Antishadowing is a
separate mechanism constrained by the moment ledger, not an arbitrary
enhancement attached to every TMD.

\subsection{Pion exchange and the \(NN\pi\) sector}

The deuteron Fock expansion contains
\begin{equation}
|D\rangle=\sqrt{Z_{NN}}|NN\rangle+|NN\pi\rangle+\cdots .
\label{eq:fock}
\end{equation}
The sourced Sullivan mechanism~\cite{MillerPion} has the longitudinal convolution
\begin{equation}
F_{a/D}^{(\pi)}(x,\kT)
=\int_x^1\frac{\dd y}{y}
\int\dd^2\bm q_T\,
f_{\pi/D}(y,\bm q_T;\Lambda',\Lambda)\,
F_{a/\pi}\left(\frac{x}{y},
\kT-\frac{x}{y}\bm q_T\right).
\label{eq:pion-convolution}
\end{equation}
The implementation uses the Miller spin-resolved splitting information and
a Vpion/JAM pion input where the corresponding channel is supported.  The
pion term is included once.  A recent LF Hamiltonian EFT calculation does
not publish the spin-resolved three-body deuteron amplitude required to
replace this boundary, so no arbitrary transfer Gaussian is invented.

\subsection{Other non-nucleonic sectors}

Delta--Delta, SRC/cluster, and hidden-color sectors have explicit
interfaces and correlated zero-centered sensitivity members.  A preferred
nonzero central addition requires a sourced flavor-, color-, helicity-, and
transverse-resolved correlator.  The present hidden-color central term is
therefore excluded rather than assigned an arbitrary probability.

\section{\texorpdfstring{\(b_1\)}{b1}, tensor normalization, and physical reductions}
\label{sec:b1}

The inclusive tensor distribution is
\begin{equation}
\delta_Tq(x)=q^0(x)-\frac12[q^{+1}(x)+q^{-1}(x)],
\label{eq:tensor-pdf}
\end{equation}
and at leading order
\begin{equation}
b_1(x,Q^2)=\frac12\sum_q e_q^2
\left[\delta_Tq(x,Q^2)+\delta_T\bar q(x,Q^2)\right].
\label{eq:b1}
\end{equation}
The project convention maps the named LL distribution by
\begin{equation}
f_{1LL}=-\frac23\,\delta_T f_1,
\label{eq:f1ll-map}
\end{equation}
with the sign adapter tested explicitly.  The HERMES comparison
\cite{HERMESb1} combines
the one-body tensor impulse term with the allowed pion contribution.  It is
a validation of the collinear LL marginal and does not determine all
tensor TMDs.

For rank-zero channels,
\begin{equation}
\int\dd^2\kT\,F_{a/D}^{(0)}(x,\kT;Q)=f_{a/D}(x;Q).
\label{eq:rankzero-marginal}
\end{equation}
For \(r>0\), the unweighted angular integral vanishes, while a weighted
moment may remain:
\begin{equation}
F^{(n)}(x)=\int\dd^2\kT
\left(\frac{k_T^2}{2M_D^2}\right)^nF(x,k_T^2).
\label{eq:transverse-moment}
\end{equation}

\section{Composition graph and uncertainty semantics}
\label{sec:composition}

The total parent is not a sum of every available table.  Each component is
typed as
\begin{equation}
\mathcal W_{\rm canonical}
=\mathcal W_{\rm baseline}
\sum_{m\in\mathcal A}\mathcal W_m^{\rm additive},
\qquad
\theta_e\in\mathcal E_e
\quad\text{for exclusive alternatives and ensembles}.
\label{eq:composition}
\end{equation}
The graph rejects two paths that represent the same phase, OAM
interference, or nuclear response.  In particular:
\begin{itemize}
  \item the CP-map nuclear responses replace legacy coefficient responses;
  \item direct axial phases and screened one-gluon phases are alternatives;
  \item spectator-inspired downstream gluon rescaling is benchmark-only;
  \item the \(NN\pi\) contribution is added once;
  \item unsupported cluster/hidden-color amplitudes are zero-centered
  alternatives.
\end{itemize}

\begin{longtable}{L{0.22\textwidth}L{0.69\textwidth}}
\caption{Meaning of uncertainty and sensitivity axes.}
\label{tab:uncertainty}\\
\toprule
Axis & Interpretation \\
\midrule
\endfirsthead
\toprule
Axis & Interpretation \\
\midrule
\endhead
PDF/TMD replica &
Member variation supplied by an external statistical or posterior ensemble:
BDSSV24 (600), JAMDiFF (968), BPV20 (500).\\
Hessian &
Correlated eigenvector propagation: CT18 (29 pairs) and MSHT20QED neutron
CSB (38 pairs).\\
Published interval hull &
All parameter corners of reported asymmetric intervals, used for Yang
\(g_{1T}\); not a confidence region without the source covariance.\\
Wave-function envelope &
AV18, CD-Bonn, NV-Ia, NV-Ib, NV-IIa, and NV-IIb dependence after
normalization and common convolution.\\
Transverse-profile envelope &
Nonperturbative \(k_T\) or \(b_T\) shape variation; it is distinct from
collinear normalization uncertainty.\\
Gauge/color scenario &
Unknown absorptive phase and f/d gluon color-link composition.\\
Nuclear mechanism &
Off-shell, shadowing, antishadowing, pion, cluster, and SRC response choices,
with channel and x-domain restrictions.\\
OAM/Fock scenario &
Alternative allowed interference content in \(L_z=0,\pm1,\pm2\) sectors.\\
CSB power counting &
Zero-centered 5\% quark and 2\% gluon amplitude envelope for TMDs lacking a
specific calculation; not a fit uncertainty.\\
Numerical error &
Quadrature, finite transform range, interpolation, and closure residuals.\\
\bottomrule
\end{longtable}

No quadrature combination is advertised as a global confidence band.
Production plots expose the named axes separately.

\section{Evidence ledger}
\label{sec:evidence}

\subsection{Quark and antiquark sectors}

\begin{longtable}{L{0.18\textwidth}L{0.31\textwidth}L{0.42\textwidth}}
\caption{Evidence status of the quark basis.  Tensor entries apply to quarks
and antiquarks with explicit flavor and proton/neutron parents.}
\label{tab:quark-evidence}\\
\toprule
Function & Central source & Interpretation and band \\
\midrule
\endfirsthead
\toprule
Function & Central source & Interpretation and band \\
\midrule
\endhead
\(f_1\) & CT18NNLO; MSHT20QED CSB &
\fit; CT18 Hessian, sourced neutron CSB, wave/profile/nuclear axes.\\
\(g_1\) & BDSSV24-NLO &
\fit; all 600 replicas and spin-dependent nuclear convolution.\\
\(h_1\) & JAMDiFF \(w\)LQCD &
\status{Fit+lattice}; 968 replicas with composed-boundary positivity
compatibility.\\
\(f_{1T}^{\perp}\) & BPV20 arTeMiDe &
\fit; 500 replicas, flavor dependence, staple reversal.\\
\(g_{1T}\) & Yang et al.\ 2024 &
\fit central; 16-corner published-interval hull and sea/OAM sensitivity.\\
\(h_{1L}^{\perp}\) & JAMDiFF \(h_1\)+WW &
\status{Fit+lattice+model}; correlated replicas and genuine-term scenario.\\
\(h_1^\perp\) & BPV20-linked hierarchy &
\model; independent Boer--Mulders operator scenario, not joint fit.\\
\(h_{1T}^{\perp}\) & positivity-bounded OAM &
\model; zero perturbative central plus signed hadronic members.\\
all LL/LT/TT &
shared S--D/OAM/Wilson parent and nuclear response &
\status{Wave-function+model+phenomenology}; common-parent prediction,
six-wave and phase/nuclear ensembles, \(b_1\) where applicable.\\
\bottomrule
\end{longtable}

\subsection{Gluon sectors}

\begin{longtable}{L{0.22\textwidth}L{0.30\textwidth}L{0.39\textwidth}}
\caption{Evidence status of the gluon basis.}
\label{tab:gluon-evidence}\\
\toprule
Function & Central source & Interpretation and band \\
\midrule
\endfirsthead
\toprule
Function & Central source & Interpretation and band \\
\midrule
\endhead
\(f_1^g\) & BSV19/NNPDF31; CT18 response &
\fit central with CT18 Hessian, wave, profile, and nuclear axes.\\
\(g_1^g\) & BDSSV24-NLO &
\fit; all 600 replicas and separate transverse-profile envelope.\\
all other U/L/T &
gluon LF overlap/OAM/Wilson parent &
\model constrained by helicity algebra, rank, f/d separation,
wave functions, positivity, and literature benchmarks.\\
all LL/LT/TT &
same shared parent with spin-1 density and nuclear maps &
\model prediction with pure-S, zero-phase, link, rank, wave, and nuclear
tests.\\
\bottomrule
\end{longtable}

This is the intended meaning of ``bringing every TMD to the \(f_1\)
level'': comparable provenance and validation discipline, not equal data
precision.

\section{Validation and numerical acceptance}
\label{sec:validation}

\subsection{Exact and controlled-limit tests}

The implemented validation matrix includes:
\begin{enumerate}
  \item correlator hermiticity, light-front parity, and gauge-link reversal;
  \item reconstruction of all allowed target and parton polarization blocks;
  \item transverse-rank covariance and physical origin limits;
  \item proton/neutron traceability and the exact charge-symmetry limit;
  \item pure-S, zero-phase, weak-phase, and real one-body limits;
  \item independent gluon f/d components and refusal of implicit hard-color
  weights;
  \item GTMD-to-TMD, TMD-to-PDF, GPD/form-factor, and \(b_1\) marginals;
  \item wave-function normalization, number, plus-momentum, and tensor
  ledgers;
  \item no-double-counting closure for nuclear and non-nucleonic mechanisms.
\end{enumerate}

\subsection{Density positivity}

For every retained joint density,
\begin{equation}
\rho=\rho^\dagger,\qquad
v^\dagger\rho v\geq0
\quad\forall\,v,
\label{eq:positivity}
\end{equation}
is tested numerically.  Positivity is imposed on the full matrix, not by
clipping individual TMDs.  At the WP12 scientific inspection,
\begin{align}
\lambda_{\min}^{qD}&=4.12612\times10^{-4},&
\lambda_{\min}^{gD}&=2.0881\times10^{-2}.
\label{eq:min-eigs}
\end{align}
No final positivity contraction was required; the common completion scale
was one.

\subsection{Closure and response diagnostics}

The resolved parent contains 38,880 quark rows and 49,680 gluon rows.
The maximum closure residual is
\begin{equation}
\epsilon_{\rm close}^{q}=0,\qquad
\epsilon_{\rm close}^{g}=1.7347\times10^{-18}.
\label{eq:closure}
\end{equation}
The maximum CP-map recomposition shifts are approximately
\(2.73\%\) and \(2.91\%\) of the local \(f_1\) reference for quarks and
gluons.  The recorded future/past residual is zero for quarks and of order
\(10^{-9}~\mathrm{GeV}^{-2}\) for gluons.

\subsection{Physical benchmarks}

Benchmarks include CT18 and BDSSV24 marginal reproduction, JAMDiFF
transversity and BPV20 Sivers member checks, the HERMES \(b_1\) comparison,
deuteron body form factors, LF current/angular-condition diagnostics,
six-wave-function comparisons, rank-aware Fourier round trips, and
literature node/hierarchy comparisons for model-dependent gluon functions.

The WP12-E acceptance report records 36/36 evidence rows and six final
gates passing.  The full repository regression after the documentation
integration is 482 passing tests.

\section{Scale scheme and the boundary before evolution}
\label{sec:evolution}

The in-house boundary uses a Collins square-root-soft subtraction
\cite{Collins,EIS}, the
delta rapidity regulator, \(\overline{\rm MS}\) UV renormalization, and the
canonical line
\begin{equation}
\zeta_i=\mu_i^2,\qquad
\mu_f=Q,\qquad \zeta_f=Q^2.
\label{eq:scale-line}
\end{equation}
The TMD equations are
\begin{align}
\frac{\dd}{\dd\ln\mu}\ln\widetilde F_a(x,b_T;\mu,\zeta)
 &=\gamma_F^a\!\left(\alpha_s,\frac{\zeta}{\mu^2}\right),\\
\frac{\dd}{\dd\ln\sqrt{\zeta}}\ln\widetilde F_a(x,b_T;\mu,\zeta)
 &=-\mathcal D_a(b_T;\mu),\\
\frac{\dd}{\dd\ln\mu}\mathcal D_a(b_T;\mu)
 &=\Gamma_{\rm cusp}^a(\alpha_s).
\label{eq:evolution}
\end{align}
At leading power, the Collins--Soper kernel depends on the quark or gluon
color representation, not on target polarization.  Polarization resides
in the boundary and matching coefficients.

The present code has heterogeneous evolution depth: BPV20 uses its native
optimal-TMD evolution; some rank-zero quark and gluon channels have the
typed in-house CSS route; several higher-rank and tensor boundaries remain
at \(Q=5~\mathrm{GeV}\) with declared profile scenarios.  Putting them on
one fixed-\(Q\) plot does not imply one complete evolution solution.

The next phase must evolve
\[
\{\text{proton-in-D},\text{ neutron-in-D},\text{ sum},
\text{ difference},\text{ nuclear correction},\text{ total}\}
\]
for every flavor, target channel, transverse rank, link/color class,
mechanism, and correlated ensemble member.  It must use the appropriate
\(J_r\) transforms, preserve closure after evolution, and refuse scheme
mismatches.

At the observable level, a complete phase-space prediction also requires
\begin{equation}
\frac{\dd\sigma}{\dd\mathcal P}
=W(q_T,Q)+Y(q_T,Q),\qquad
Y=\sigma_{\rm FO}-[W]_{\rm expanded\ to\ FO},
\label{eq:wy}
\end{equation}
so the nonperturbative TMD boundary is not forced to mimic the high-\(q_T\)
fixed-order tail.  A general process-specific \(Y\) term remains open.

\section{Implemented software map and reproducibility}
\label{sec:software}

\begin{longtable}{L{0.25\textwidth}L{0.65\textwidth}}
\caption{Principal computational objects.}
\label{tab:software}\\
\toprule
Layer & Representative implementation \\
\midrule
\endfirsthead
\toprule
Layer & Representative implementation \\
\midrule
\endhead
Spin and transverse representation &
\texttt{spin.py}, \texttt{transverse\_tensors.py},
\texttt{quark\_correlator.py}, \texttt{gluon\_correlator.py},
\texttt{registry.py}.\\
Nucleon quark inputs &
\texttt{nucleon\_inputs.py}, \texttt{transversity.py},
\texttt{bpv20\_sivers.py}, \texttt{worm\_gear\_inputs.py},
\texttt{boer\_mulders.py}.\\
Nucleon gluon and phases &
\texttt{nucleon\_gluon\_inputs.py}, \texttt{gluon\_lfwf\_todd.py},
\texttt{gluon\_todd.py}, \texttt{gauge\_link\_phase.py}.\\
Deuteron wave functions &
\texttt{wavefunctions/av18.py}, \texttt{cd\_bonn.py},
\texttt{norfolk.py}, \texttt{light\_front.py}.\\
Nuclear mechanisms &
\texttt{gtmd\_convolution.py}, \texttt{operator\_nuclear\_response.py},
\texttt{diffractive\_shadowing.py}, \texttt{pion\_exchange.py},
\texttt{nonnucleonic\_transverse.py}.\\
Composition and resolved output &
\texttt{canonical\_composition.py},
\texttt{canonical\_parent\_enrichment.py},
\texttt{resolved\_nuclear\_parent.py}.\\
Scheme and transforms &
\texttt{tmd\_scheme.py}, \texttt{canonical\_scheme.py},
\texttt{quark\_tmd\_matching.py}, \texttt{gluon\_tmd\_matching.py},
\texttt{tmd\_evolution.py}, \texttt{fourier.py}.\\
Validation &
\texttt{joint\_positivity.py}, \texttt{moment\_ledger.py},
\texttt{controlled\_limits.py}, and the generated WP11/WP12 reports.\\
\bottomrule
\end{longtable}

The authoritative numerical parents are
\begin{itemize}
  \item \texttt{outputs/parent\_tmds/wp12\_resolved\_quark\_parent.csv};
  \item \texttt{outputs/parent\_tmds/wp12\_resolved\_gluon\_parent.csv};
  \item their \texttt{.correlators.csv} companions;
  \item \texttt{outputs/parent\_tmds/wp12\_canonical\_composed\_quark.csv};
  \item \texttt{outputs/parent\_tmds/wp12\_canonical\_composed\_gluon.csv}.
\end{itemize}
The acceptance authority is
\texttt{outputs/validation/wp12e\_acceptance.json} together with the
resolved-parent and evidence-parity reports.

\section{Requirements for a genuinely predictive next-level model}
\label{sec:next-level}

\subsection{Change of model class}

The present construction is a symmetry-complete and nuclear-dressed
phenomenological synthesis.  Its common correlator architecture enforces
consistency, but many of the scalar amplitudes entering that correlator come
from different fits or from separately normalized model sectors.  This is
not equivalent to deriving the full correlator from one microscopic quantum
state.

A genuinely predictive next-level model requires a change in the direction
of inference.  The present construction follows
\begin{equation}
\{\text{PDF/TMD fits, nuclear models, phase models}\}
\longrightarrow
\Phi_{a/D}^{\rm assembled}
\longrightarrow
\{F_A\}.
\label{eq:assembled-direction}
\end{equation}
The next-level construction must instead follow
\begin{equation}
\{H_{\rm LF}^{\rm ren},\bm\theta_{\rm micro}\}
\longrightarrow
\{|\Psi_N\rangle,|\Psi_D\rangle,\mathcal U_\eta\}
\longrightarrow
W_{a/N},W_{a/D}
\longrightarrow
\{F_A,\mathcal O_{\rm process}\},
\label{eq:generative-direction}
\end{equation}
where one renormalized microscopic model and a limited shared parameter
vector \(\bm\theta_{\rm micro}\) generate all flavors, polarizations, OAM
interferences, gauge-link phases, nuclear components, and observables.
External PDFs, TMDs, lattice matrix elements, form factors, and \(b_1\) then
calibrate or test the same parameter vector; they no longer supply unrelated
named TMDs.

Even this stronger construction would be a controlled effective-QCD
prediction, not an exact solution of continuum QCD.  Its predictive content
would come from using substantially fewer shared microscopic degrees of
freedom than the number of observables predicted, and from successfully
predicting quantities that were not used for calibration.

\subsection{Microscopic light-front Hamiltonian and regulated Hilbert space}

The first required object is a gauge-consistent, renormalized light-front
Hamiltonian eigenproblem,
\begin{equation}
H_{\rm LF}^{\rm ren}(\Lambda_{\rm UV},\Lambda_{\rm IR};
\bm\theta_{\rm QCD},\bm\theta_{\rm EFT})
|\Psi_h(P^+,\bm P_T,\Lambda)\rangle
=M_h^2|\Psi_h(P^+,\bm P_T,\Lambda)\rangle .
\label{eq:lf-eigenproblem}
\end{equation}
The regulator, counterterms, basis truncation, and renormalization conditions
must be part of the model definition.  A numerical diagonalization without
a cutoff study or sector-dependent renormalization is not a predictive
Hamiltonian model.

For the nucleon, the minimum useful Fock expansion is
\begin{align}
|N,\lambda\rangle
&=\psi_{qqq}^{\lambda}|qqq\rangle
 +\psi_{qqqg}^{\lambda}|qqqg\rangle
 +\psi_{qqqq\bar q}^{\lambda}|qqqq\bar q\rangle
 +\cdots ,
\label{eq:nucleon-fock}
\end{align}
with explicit \(u,d,\bar u,\bar d\) and gluon degrees of freedom, color
singlet constraints, parton helicities, and orbital sectors sufficient to
support the transverse ranks in \cref{tab:quark-basis,tab:gluon-basis}.
At minimum, the basis must contain the \(L_z=0,\pm1,\pm2\) interferences
needed by Sivers, Boer--Mulders, worm-gear, pretzelosity, linear-gluon, and
spin-1 tensor projections.  A valence-only \(|qqq\rangle\) state cannot
predict antiquark TMDs, genuine quark--gluon correlations, or the gluon
sector.

The Hamiltonian or effective interaction must specify:
\begin{enumerate}
  \item quark masses and kinetic terms in a declared scheme;
  \item confining interactions and their rotational-symmetry restoration;
  \item quark--gluon vertices and instantaneous light-front interactions;
  \item chiral dynamics capable of producing the pion cloud and sea-flavor
  asymmetry;
  \item counterterms and sector-dependent renormalization conditions;
  \item zero-mode treatment and the consequences of Fock-space truncation;
  \item a convergence sequence in longitudinal resolution, transverse basis,
  Fock sectors, and ultraviolet cutoff.
\end{enumerate}

Candidate realizations include basis light-front quantization, a
Hamiltonian effective-field-theory construction, or a continuum
Bethe--Salpeter/Dyson--Schwinger amplitude mapped consistently to the light
front.  The project should select one primary generative route.  Combining
wave functions from mutually unrelated microscopic models would reproduce
the present patchwork at a deeper level.

\subsection{Nucleon GTMDs from common wave-function overlaps}

For T-even distributions, the renormalized nucleon GTMD must be calculated
as an overlap of the same Fock amplitudes.  Schematically,
\begin{align}
W_{a/N,\lambda'\lambda}^{[\Gamma]}
(x,\kT,\DT)
&=\sum_{n}\sum_{j\in a}
\int[\dd x]_n[\dd^2\bm k_T]_n\,
\delta(x-x_j)\delta^{(2)}(\kT-\bm k_{Tj})
\nonumber\\
&\quad\times
\psi_{n,\lambda'}^\ast(\{x_i,\bm k'_{Ti},\lambda_i,c_i\})
\mathcal P_{j}^{[\Gamma]}
\psi_{n,\lambda}(\{x_i,\bm k_{Ti},\lambda_i,c_i\}),
\label{eq:microscopic-overlap}
\end{align}
where the initial/final transverse momenta implement one common recoil
convention and \(\mathcal P_j^{[\Gamma]}\) acts on the active parton.
Quark, antiquark, and gluon correlators must be derived from this same
state.  Their normalizations are then linked by baryon number, momentum,
helicity, and angular-momentum sum rules rather than independently chosen.

The overlap implementation must demonstrate that
\begin{equation}
\left.
\begin{array}{c}
W(x,\kT,\DT)\\[1mm]
\displaystyle\int\dd^2\kT\,W(x,\kT,\DT)\\[1mm]
W(x,\kT,\bm0)\\[1mm]
\displaystyle\int\dd x\,\dd^2\kT\,W
\end{array}
\right\}
\longrightarrow
\left\{
\begin{array}{c}
\text{GTMD}\\
\text{GPD}\\
\text{TMD}\\
\text{local current or EMT form factor}
\end{array}
\right.
\label{eq:microscopic-reductions}
\end{equation}
with one regulator and one state normalization.  A TMD-specific Gaussian
may be used only as a quantified numerical regulator or controlled
approximation, not as an independent physical width for every function.

The microscopic state must also reproduce OAM and spin decompositions.  For
example, in a declared gauge and scheme,
\begin{equation}
\frac12
=\frac12\Delta\Sigma(\mu)
+L_q(\mu)+\Delta G(\mu)+L_g(\mu),
\label{eq:spin-sum}
\end{equation}
and the canonical/kinetic OAM distinction must be related to the gauge-link
definition used by the GTMD.  The sign pattern of flavor-resolved Sivers or
OAM-sensitive functions must emerge from these amplitudes, not be imposed
as a plot-level coefficient.

\subsection{Dynamical Wilson lines and naive-T-odd functions}

Real light-front overlaps alone cannot predict naive-T-odd TMDs.  The
microscopic model must calculate, or systematically approximate, the
process-dependent Wilson line
\begin{equation}
\mathcal U_\eta[0,z]
=\mathcal P\exp\left[
ig\int_{\mathcal C_\eta(0,z)}
\dd\xi^\mu A_\mu^a(\xi)t^a
\right],
\label{eq:wilson-line}
\end{equation}
using the same gluon fields and color couplings that appear in
\cref{eq:lf-eigenproblem}.  An expansion such as
\begin{equation}
\mathcal U_\eta
=1+ig\int_{\mathcal C_\eta}\dd\xi\cdot A
(ig)^2\int_{\mathcal C_\eta}\!\!\int_{\mathcal C_\eta}
\mathcal P[A(\xi)A(\zeta)]+\cdots
\label{eq:wilson-expansion}
\end{equation}
must have a declared truncation, rapidity regulator, color representation,
and infrared treatment.

At minimum, the calculation must:
\begin{itemize}
  \item generate the absorptive phase from explicit interference with
  intermediate states, rather than assign an arbitrary imaginary constant;
  \item recover exact future/past sign reversal where factorization requires
  it;
  \item retain distinct quark flavor and gluon \(f^{abc}\)- and
  \(d^{abc}\)-type color structures;
  \item reproduce a perturbative one-gluon limit and demonstrate stability
  when higher eikonal orders are included;
  \item identify processes for which the assumed staple factorization is
  established, conditional, or broken by Glauber/color entanglement.
\end{itemize}

The quark Sivers, Boer--Mulders, and all gluon T-odd functions would then be
different projections of common rescattering amplitudes.  Their relative
signs and magnitudes would become correlated predictions.  A universal phase
multiplied onto unrelated operators is specifically disallowed.

\subsection{Microscopic spin-1 nuclear state}

The deuteron cannot be generated by convolving a microscopic nucleon with an
unrelated phenomenological nuclear response.  The nuclear state must be
derived in a Hilbert space that resolves both nucleonic and non-nucleonic
sectors:
\begin{align}
|D,\Lambda\rangle
&=\Psi_{NN}^{\Lambda}|NN\rangle
 +\Psi_{NN\pi}^{\Lambda}|NN\pi\rangle
 +\Psi_{\Delta\Delta}^{\Lambda}|\Delta\Delta\rangle
 +\Psi_{6q}^{\Lambda}|6q\rangle
 +\cdots .
\label{eq:deuteron-fock-next}
\end{align}
The amplitudes must share one normalization,
\begin{equation}
1=Z_{NN}+Z_{NN\pi}+Z_{\Delta\Delta}+Z_{6q}+\cdots ,
\qquad
1=\sum_\alpha\int[\dd\Gamma_\alpha]\,
\sum_i x_i|\Psi_\alpha|^2,
\label{eq:fock-normalization}
\end{equation}
so baryon number and plus momentum close without after-the-fact
renormalization.

The \(NN\) sector must preserve the successful AV18/CD-Bonn/Norfolk
long-distance S/D phenomenology or derive an equally accurate interaction,
including consistent one- and two-body currents.  The higher sectors must
carry explicit spin, isospin, color, parton flavor, transverse momentum,
and off-forward transfer.  A scalar pion distribution or a six-quark
probability is insufficient to predict tensor TMDs.

The deuteron GTMD should be evaluated as a matrix element in the full state,
\begin{equation}
W_{a/D}^{[\Gamma]}
=\sum_{\alpha,\beta}
\langle\Psi_D^\alpha|
\widehat{\mathcal O}_{a}^{[\Gamma]}[\mathcal U_\eta]
|\Psi_D^\beta\rangle ,
\label{eq:full-deuteron-matrix}
\end{equation}
including diagonal contributions and physically allowed interference
between Fock sectors.  Impulse, exchange-current, shadowing, and
non-nucleonic terminology should then label controlled reductions of
\cref{eq:full-deuteron-matrix}, not independently normalized additions.

At small x, coherent multiple scattering must be derived from the same
nucleon diffractive amplitudes and nuclear state.  Polarized and tensor
shadowing cannot be copied from the unpolarized response; their helicity
amplitudes must be calculated.  At intermediate and large x, off-shell,
Fermi-motion, SRC, and EMC-like behavior must emerge from the joint
spectral/partonic amplitude with a tested separation of scales.

\subsection{Matching, renormalization, and evolution}

A low-resolution Hamiltonian wave function is not by itself a TMD in a
specified QCD scheme.  A complete prediction requires a matching map
\begin{equation}
\widetilde W^{\rm QCD}_{a/h}
(x,b_T;\mu,\zeta)
=\mathcal Z_{a}^{\rm LF\to TMD}
(b_T;\Lambda_{\rm LF},\mu,\zeta)
\otimes
\widetilde W^{\rm LF}_{a/h}
(x,b_T;\Lambda_{\rm LF})
+\mathcal O\!\left(\frac{\Lambda_{\rm LF}}{Q}\right),
\label{eq:lf-matching}
\end{equation}
including soft subtraction, rapidity renormalization, and operator mixing.
The matching must identify which Hamiltonian regulator dependence cancels
against \(\mathcal Z^{\rm LF\to TMD}\).

The next-level model must use one consistent route for:
\begin{enumerate}
  \item nonsinglet, singlet, and quark--gluon collinear evolution;
  \item quark and gluon Collins--Soper kernels;
  \item transverse ranks \(r=0,1,2,\ldots\);
  \item T-even and T-odd operators in their correct gauge-link classes;
  \item threshold matching and heavy flavors when the Q range requires them;
  \item process-level \(W+Y\) matching.
\end{enumerate}
Fit-native evolved inputs may remain as external validation data, but they
must not replace the evolution of the microscopic state.

\subsection{Calibration without rebuilding a patchwork}

The microscopic parameter vector must be small, shared, and identifiable.
The likelihood should be defined at the observable or well-defined matrix
element level:
\begin{equation}
p(\bm\theta_{\rm micro}|\mathcal D)
\propto
p(\mathcal D|\bm\theta_{\rm micro})
p_{\rm EFT}(\bm\theta_{\rm micro})
p_{\rm ren}(\bm\theta_{\rm micro}),
\label{eq:micro-posterior}
\end{equation}
with correlated experimental, lattice, truncation, regulator, and numerical
covariances.  Suitable calibration data include:
\begin{itemize}
  \item hadron masses, charges, magnetic moments, and axial charges;
  \item nucleon electromagnetic, axial, and energy--momentum form factors;
  \item selected unpolarized/helicity PDFs and low Mellin moments;
  \item lattice quasi-/pseudo-distributions and GTMD-related moments where
  scheme matching is controlled;
  \item deuteron binding, asymptotic S/D normalization, elastic form factors,
  quadrupole and magnetic moments;
  \item a restricted subset of \(b_1\), SIDIS, and T-odd data sufficient to
  constrain common phases.
\end{itemize}

The calibration must not assign an independent free normalization to each
TMD.  A parameter associated with a quark--gluon vertex, OAM mixing, or
eikonal kernel must affect every correlator projection to which that
microscopic interaction contributes.  Identifiability should be assessed
through posterior correlations, Fisher information or profile likelihoods,
and closure on withheld observables.

The fit must be divided into calibration and prediction sets before final
optimization.  A quantity counts as a prediction only if its normalization
and shape were not directly fitted.  For example, using Sivers data to fix a
common rescattering kernel and then predicting Boer--Mulders, tensor T-odd,
and gluon color-sector observables is a meaningful test; fitting each of
those independently is not.

\subsection{Uncertainty budget for a generative model}

The next-level uncertainty is not a band assembled from unrelated models.
It must propagate a common ensemble of microscopic states.  For any
observable \(\mathcal O\),
\begin{equation}
\operatorname{Cov}(\mathcal O_i,\mathcal O_j)
=\left\langle
\big[\mathcal O_i(\bm\theta,\nu)-\overline{\mathcal O}_i\big]
\big[\mathcal O_j(\bm\theta,\nu)-\overline{\mathcal O}_j\big]
\right\rangle_{\bm\theta,\nu},
\label{eq:predictive-covariance}
\end{equation}
where \(\bm\theta\) denotes calibrated microscopic parameters and \(\nu\)
collects controlled truncation, regulator, matching, and numerical choices.
The covariance must retain cross-flavor, cross-TMD, quark--gluon,
proton--neutron, and cross-observable correlations.

The error budget must separately report:
\begin{enumerate}
  \item statistical/calibration uncertainty;
  \item Hamiltonian EFT truncation;
  \item Fock-sector and basis truncation;
  \item ultraviolet/infrared regulator and renormalization uncertainty;
  \item Wilson-line/eikonal truncation;
  \item perturbative matching and scale variation;
  \item TMD evolution and large-\(b_T\) uncertainty;
  \item nuclear many-body truncation;
  \item numerical discretization and transform error.
\end{enumerate}
Convergence in one category cannot be used to hide another.  In particular,
positivity is a constraint on allowed states, not an uncertainty estimate.

\subsection{Validation ladder and falsifiability}

The model should be validated in increasing dynamical depth:
\begin{enumerate}
  \item \textbf{Algebraic layer:} color singlets, Poincare/light-front
  kinematics, Hermiticity, parity, time reversal with link reversal,
  permutation symmetry, and exact spin-1 projection closure.
  \item \textbf{State layer:} normalization, regulator independence within
  truncation errors, mass spectrum, rotational multiplets, and convergence
  of Fock probabilities and OAM components.
  \item \textbf{Local-current layer:} charges, magnetic moments, axial
  charges, energy--momentum moments, and Ward identities using currents
  consistent with the Hamiltonian.
  \item \textbf{Marginal layer:} simultaneous TMD/GPD/PDF/form-factor
  reductions from the same GTMD.
  \item \textbf{Nucleon layer:} withheld flavor-resolved PDFs, TMD moments,
  lattice observables, and process sign changes.
  \item \textbf{Deuteron layer:} binding, elastic observables, \(b_1\),
  tagged structure, tensor moments, and controlled proton/neutron limits.
  \item \textbf{Process layer:} SIDIS, DY, and gluon-sensitive predictions
  with explicit hard factors, fragmentation, color weights, and \(W+Y\).
\end{enumerate}

Failure of a withheld observable must update or reject the microscopic
model.  It must not be repaired by adding a TMD-specific coefficient unless
that coefficient corresponds to a missing operator added consistently to
the Hamiltonian or current.

\subsection{Required software architecture and deliverables}

The present modular project can host the new model, but the central state
must change.  Required new computational objects are:
\begin{enumerate}
  \item a versioned \texttt{MicroscopicHamiltonian} carrying regulator,
  counterterm, basis, and renormalization metadata;
  \item a sparse/distributed \texttt{FockState} representation with parton
  flavor, color, helicity, momentum, OAM, and sector identity;
  \item eigensolvers and convergence manifests for nucleon and deuteron
  states;
  \item common quark/gluon GTMD overlap evaluators at nonzero \(\DT\);
  \item a Wilson-line engine with perturbative/eikonal order and link/color
  labels;
  \item a matching adapter from the Hamiltonian regulator to the declared
  soft-subtracted TMD scheme;
  \item an ensemble store that preserves microscopic member identity through
  every observable;
  \item automatic-differentiation or adjoint sensitivities for calibration;
  \item observable-level likelihoods and withheld-validation reports.
\end{enumerate}

PennyLane or another quantum simulator would be justified only if it
represents a defined truncated Fock Hilbert space, implements the same
Hamiltonian matrix elements, and is benchmarked against exact diagonalization
on smaller sectors.  A quantum circuit that merely reproduces Clebsch--Gordan
coupling does not add predictive physics.

\subsection{Acceptance criteria}

The next-level model may be called genuinely predictive only when all of the
following are satisfied:
\begin{enumerate}
  \item One declared microscopic Hamiltonian/state ensemble generates
  quark, antiquark, and gluon correlators; fitted named-TMD curves are
  comparison data, not central inputs.
  \item The nucleon state contains the Fock sectors needed for sea and gluon
  observables, with demonstrated basis and sector convergence.
  \item Proton and neutron flavor differences, helicity, transversity, OAM,
  and spin--orbit correlations emerge from the state.
  \item T-even GTMDs and all their TMD/GPD/PDF/form-factor reductions close
  from the same overlap.
  \item T-odd phases are generated by an explicit Wilson-line/rescattering
  calculation with correct process and f/d color dependence.
  \item The spin-1 state contains normalized \(NN\) and supported
  non-nucleonic sectors with explicit tensor polarization and interference.
  \item Currents and operators are matched consistently to the Hamiltonian
  and satisfy Ward identities and known local limits.
  \item One QCD TMD scheme and rank-aware evolution route carries the full
  state to multiple Q values.
  \item A shared microscopic posterior predicts withheld nucleon and
  deuteron observables with correlated uncertainties.
  \item At least one nontrivial family not used in calibration---for example
  tensor TMDs, Boer--Mulders from a Sivers-calibrated kernel, or gluon
  f/d-sensitive observables---is successfully predicted.
  \item Removing any presently phenomenological TMD input does not make the
  central microscopic calculation undefined.
  \item All symmetry, positivity, marginal, sum-rule, regulator,
  truncation, numerical, and no-double-counting gates pass.
\end{enumerate}

\begin{table}[ht]
\centering
\caption{Distinction between the present synthesis and the required
next-level generative model.}
\label{tab:model-class-comparison}
\small
\begin{tabularx}{\textwidth}{L{0.22\textwidth}Y Y}
\toprule
Question & Present model & Genuinely predictive model \\
\midrule
Source of TMD amplitudes &
Different fits and explicit model sectors &
Common microscopic nucleon/deuteron state \\
Flavor and spin &
Retained and constrained after input assembly &
Generated by Hamiltonian eigenstates \\
T-odd phases &
Fit input or screened/eikonal scenarios &
Wilson-line dynamics using the model gluon field \\
Tensor nuclear physics &
Realistic wave functions plus operator response maps &
Full spin-resolved nuclear Fock-state matrix element \\
Non-nucleonic sectors &
Sourced additions or zero-centered interfaces &
Normalized amplitudes in the same eigenstate \\
Uncertainty &
Named heterogeneous axes &
Correlated microscopic posterior plus truncation errors \\
Evolution &
Heterogeneous fixed-\(Q\) boundary routes &
One matched, rank-aware QCD scheme \\
Meaning of prediction &
Constrained extrapolation or model scenario &
Withheld observable from shared calibrated parameters \\
\bottomrule
\end{tabularx}
\end{table}

This program is substantially more demanding than completing evolution of
the current boundary.  Evolution remains useful for the present
phenomenological model, but it must not be presented as the step that turns
that model into a microscopic prediction.  The Hamiltonian/Fock-state and
Wilson-line program above is the actual transition in model class.

\section{Completion audit and remaining physics}
\label{sec:audit}

\subsection{Completed at the declared boundary}

\begin{itemize}
  \item All 18 quark/antiquark and 18 gluon forward projections.
  \item Explicit \(u,d,\bar u,\bar d\), proton, and neutron parents.
  \item U/L/T/LL/LT/TT target structure and retained parton helicity.
  \item Independent gluon f/d gauge-link/color components.
  \item A shared GTMD/correlator parent and tested projections.
  \item Six realistic deuteron wave functions with S/D and SS/SD/DS/DD
  interference.
  \item Spin, OAM, spin--orbit, and absorptive Wilson-line mechanisms.
  \item Binding/Fermi motion, off-shell, shadowing, antishadowing, pion, and
  controlled non-nucleonic interfaces.
  \item Named fit, Hessian, interval, wave, profile, phase, nuclear, OAM,
  CSB, and numerical axes.
  \item Symmetry, positivity, support, marginal, moment, closure, controlled
  limit, and benchmark validation.
\end{itemize}

\subsection{Model-dominated but honestly represented}

Boer--Mulders normalization, pretzelosity, genuine WW breaking, most
tensor-polarized TMDs, polarized/tensor shadowing, absolute non-Sivers
gluon T-odd magnitudes, process-specific f/d mixtures, and intrinsic
non-nucleonic transverse/color structure remain model dominated.  They are
not omitted, but neither are their bands presented as fitted confidence
regions.

\subsection{Open requirements after this boundary}

\begin{enumerate}
  \item Complete rank-aware, common-scheme, multi-\(Q\) evolution with
  resolved-component and correlated-ensemble closure.
  \item Process-specific \(W+Y\) matching and hard/color factors.
  \item A global inference program combining tensor DIS, SIDIS, DY,
  lattice constraints, and gluon-sensitive observables.
  \item Full nonzero-skewness GTMD phenomenology and polynomiality tests.
  \item Improved polarized/tensor diffractive inputs and a spin-resolved
  transverse \(NN\pi\) amplitude.
  \item Direct fit or lattice information for tensor and gluon TMDs that
  currently rely on shared-parent model predictions.
\end{enumerate}

\section{Conclusion}

The original GTMD-first proposal supplied the right organizing principle,
but the route to an adequate numerical boundary required two substantive
corrections.  Exact zeros in a real scalar one-body parent were recognized
as missing-amplitude limits, and a later smooth reduced-amplitude atlas was
recognized as a bookkeeping closure rather than a derivation from the
light-front parent.  The present model restores the parent chain, preserves
flavor and constituent identity, represents vector and tensor spin-1
structure, introduces physically distinct OAM and gauge-link amplitudes,
and composes nuclear and non-nucleonic mechanisms under an explicit
no-double-counting contract.

The strongest sectors inherit modern unpolarized, helicity, transversity,
and Sivers ensembles.  Less constrained tensor and gluon sectors are
predictions of the same spin/OAM/Wilson parent with explicit sensitivity
axes.  The final deuteron may approach isoscalar equality, but this emerges
from unequal proton and neutron flavor parents rather than being built in
as a flavor-independent ansatz.  All declared projections now meet the
same standard of provenance, replacement interfaces, nuclear propagation,
and validation, although not the same experimental precision.

This is a scientifically defensible boundary on which to build complete
evolution.  Evolution must transport, rather than erase, the flavor, spin,
rank, color/link, constituent, mechanism, and uncertainty structure
documented here.

\appendix
\section{Machine-readable acceptance values}

\begin{tabularx}{\textwidth}{L{0.48\textwidth}Y}
\toprule
Quantity & Accepted value \\
\midrule
Declared boundary &
\(Q=5~\mathrm{GeV}\),
\(x_N=\{0.02,0.05,0.10,0.20,0.40\}\) \\
Basis & 18 quark/antiquark projections per flavor; 18 gluon projections \\
Resolved parent size & 38,880 quark rows; 49,680 gluon rows \\
Evidence parity & 36/36 rows pass \\
Minimum density eigenvalues &
\(4.12612\times10^{-4}\) (quark),
\(2.0881\times10^{-2}\) (gluon) \\
Maximum parent closure residual &
\(0\) (quark), \(1.7347\times10^{-18}\) (gluon) \\
Maximum CP recomposition shift &
\(2.73\%\) (quark), \(2.91\%\) (gluon), relative to local \(f_1\) \\
Full regression & 482 tests pass \\
\bottomrule
\end{tabularx}

\section{Reproduction commands}

\begin{verbatim}
# Compile this manuscript.
TECTONIC_CACHE_DIR=/private/tmp/deuteron-tectonic-cache \
  /Users/dustin/work/DeuteronWigner/.conda-latex/bin/tectonic \
  -X compile references/model_construction_note.tex \
  --outdir output/pdf

# Run the complete repository validation.
PYTHONPATH=src MPLCONFIGDIR=/private/tmp/deuteron-mpl \
  /Users/dustin/miniforge3/bin/python3.9 -m pytest -q
\end{verbatim}

\small
\begin{thebibliography}{99}

\bibitem{Hoodbhoy1989}
P.~Hoodbhoy, R.~L.~Jaffe, and A.~Manohar,
``Novel effects in deep inelastic scattering from spin-one hadrons,''
Nucl.\ Phys.\ B \textbf{312}, 571 (1989).

\bibitem{HERMESb1}
A.~Airapetian et al.\ (HERMES Collaboration),
``First measurement of the tensor structure function \(b_1\) of the
deuteron,'' Phys.\ Rev.\ Lett.\ \textbf{95}, 242001 (2005),
arXiv:hep-ex/0506018.

\bibitem{Cosyn2017}
W.~Cosyn, Y.-B.~Dong, S.~Kumano, and M.~Sargsian,
``Tensor-polarized structure function \(b_1\) in the standard convolution
description of the deuteron,'' Phys.\ Rev.\ D \textbf{95}, 074036 (2017),
arXiv:1702.05337.

\bibitem{BacchettaMulders2000}
A.~Bacchetta and P.~J.~Mulders,
``Deep inelastic leptoproduction of spin-one hadrons,''
Phys.\ Rev.\ D \textbf{62}, 114004 (2000), arXiv:hep-ph/0007120.

\bibitem{KumanoSong2021}
S.~Kumano and Q.-T.~Song,
``Transverse-momentum-dependent parton distribution functions for spin-1
hadrons,'' Phys.\ Rev.\ D \textbf{103}, 014025 (2021),
arXiv:2011.08583.

\bibitem{Spin1Quark2016}
S.~Kumano and Q.-T.~Song,
``TMDs for spin-1 hadrons,'' arXiv:1612.06585.

\bibitem{Boer2016}
D.~Boer, S.~Cotogno, T.~van Daal, P.~J.~Mulders, A.~Signori, and
Y.-J.~Zhou, ``Gluon and Wilson loop TMDs for hadrons of spin
\(\leq1\),'' JHEP \textbf{10}, 013 (2016), arXiv:1607.01654.

\bibitem{Cotogno2017}
S.~Cotogno, T.~van Daal, and P.~J.~Mulders,
``Positivity bounds on gluon TMDs for hadrons of spin \(\leq1\),''
JHEP \textbf{11}, 185 (2017), arXiv:1709.07827.

\bibitem{Belitsky2004}
A.~V.~Belitsky, X.~Ji, and F.~Yuan,
``Quark imaging in the proton via quantum phase-space distributions,''
Phys.\ Rev.\ D \textbf{69}, 074014 (2004), arXiv:hep-ph/0307383.

\bibitem{Lorce2011}
C.~Lorc\'e and B.~Pasquini,
``Quark Wigner distributions and orbital angular momentum,''
Phys.\ Rev.\ D \textbf{84}, 014015 (2011), arXiv:1106.0139.

\bibitem{AV18}
R.~B.~Wiringa, V.~G.~J.~Stoks, and R.~Schiavilla,
``Accurate nucleon-nucleon potential with charge-independence breaking,''
Phys.\ Rev.\ C \textbf{51}, 38 (1995), arXiv:nucl-th/9408016.

\bibitem{CDBonn}
R.~Machleidt,
``The high-precision, charge-dependent Bonn nucleon-nucleon potential,''
Phys.\ Rev.\ C \textbf{63}, 024001 (2001), arXiv:nucl-th/0006014.

\bibitem{Neff2016}
T.~Neff and H.~Feldmeier,
``The Wigner function and short-range correlations in the deuteron,''
arXiv:1610.04066.

\bibitem{BPV20}
M.~Bury, A.~Prokudin, and A.~Vladimirov,
``Extraction of the Sivers function from SIDIS, Drell--Yan, and \(W^\pm/Z\)
data at N3LO,'' JHEP \textbf{05}, 151 (2021), arXiv:2103.03270.

\bibitem{CT18}
T.-J.~Hou et al.,
``New CTEQ global analysis of quantum chromodynamics with high-precision
data from the LHC,'' Phys.\ Rev.\ D \textbf{103}, 014013 (2021),
arXiv:1912.10053.

\bibitem{BDSSV24}
D.~de Florian, R.~Sassot, M.~Stratmann, and W.~Vogelsang,
updated BDSSV polarized parton-distribution analysis and public
BDSSV24-NLO replica set (2024); project input metadata records the released
LHAPDF member identities.

\bibitem{JAMDiff}
C.~Cocuzza et al.,
``Transversity distributions with lattice-QCD constraints,''
arXiv:2306.12998.

\bibitem{Yang2024}
X.~Yang et al.,
phenomenological world-SIDIS extraction of the transverse worm-gear
distribution \(g_{1T}\), arXiv:2403.12795, especially Eq.~(46) and
Table~IV.

\bibitem{Pretzel}
Y.~Chai, K.-B.~Chen, and J.-P.~Ma,
``A note on pretzelosity TMD parton distribution,''
arXiv:1808.10560.

\bibitem{BSV19}
I.~Scimemi and A.~Vladimirov,
``Non-perturbative structure of semi-inclusive deep-inelastic and
Drell--Yan scattering at small transverse momentum,''
JHEP \textbf{06}, 137 (2020), arXiv:1912.06532; see also the BSV19
arTeMiDe input documented in the project scheme manifest.

\bibitem{MSHTQED}
T.~Cridge, L.~A.~Harland-Lang, A.~D.~Martin, and R.~S.~Thorne,
``QED parton distribution functions in the MSHT20 fit,''
Eur.\ Phys.\ J.\ C \textbf{82}, 90 (2022), arXiv:2111.05357.

\bibitem{FGS}
L.~Frankfurt, V.~Guzey, and M.~Strikman,
``Leading twist nuclear shadowing phenomena in hard processes with
nuclei,'' Phys.\ Rept.\ \textbf{512}, 255 (2012),
arXiv:1106.2091; see also arXiv:hep-ph/0601123 for the deuteron expression.

\bibitem{H1DPDF}
A.~Aktas et al.\ (H1 Collaboration),
``Dijet cross sections and parton densities in diffractive DIS at HERA,''
JHEP \textbf{10}, 042 (2007), arXiv:0708.3217.

\bibitem{MillerPion}
G.~A.~Miller,
``Pionic and hidden-color, six-quark contributions to the deuteron
\(b_1\) structure function,'' Phys.\ Rev.\ C \textbf{89}, 045203 (2014),
arXiv:1311.4561.

\bibitem{Collins}
J.~Collins, \emph{Foundations of Perturbative QCD},
Cambridge University Press (2011).

\bibitem{EIS}
M.~G.~Echevarria, A.~Idilbi, and I.~Scimemi,
``Factorization theorem for Drell--Yan at low \(q_T\),''
JHEP \textbf{07}, 002 (2012), arXiv:1111.4996.

\bibitem{LFCao2026}
X.~Cao et al.,
``Non-perturbative flavor asymmetry in the nucleon and deuteron: the
light-front Hamiltonian effective field theory approach,''
arXiv:2601.13567 (2026).

\end{thebibliography}

\end{document}
